% ===================================================================
%  नकसा नं. ११ — सहर आ ओकर इमारत। --- आग।
%
%  Traduction, faite en 2026, en bhojpouri d'aujourd'hui, sur l'ido
%  de texto/io. Regles de la colonne : voir 10-nakasa-01.tex. DEUX
%  parties, deux numerotations qui repartent de 1 (t11-c1-*,
%  t11-c2-*). 29 blocs, 89 renvois, dix-neuf paires a retourner. DEUX
%  blocs — t11-c1-12-1 et t11-c2-03-1 — tiennent chacun DEUX alineas
%  sous une seule cle : ne pas en ouvrir de seconde.
%
%  « अस्पताल » L'HOPITAL (39) EST LE QUATRIEME MOT PORTUGAIS DE LA
%  COLONNE, ET LE COMPTE EST MAINTENANT CONCLUANT. Le releve tenait
%  मेज la table (tableau 1), पादरी le cure (tableau 3), मिस्त्री le
%  macon (tableau 5) ; voici अस्पताल, de hospital. LES QUATRE
%  FIGURAIENT DEJA DANS LA LISTE DES ONZE MOTS PORTUGAIS DE LA
%  COLONNE GUJARATIE — મેજ, પાદરી, મિસ્ત્રી, ઇસ્પિતાલ. Or le Gujarat
%  est la cote ou les Portugais ont debarque et le pays bhojpouri est
%  a mille cinq cents kilometres de la. Quatre mots sur quatre
%  franchissent la distance : ce n'est donc pas une couche cotiere,
%  c'est une couche INDIENNE, et c'est la reponse a la question que
%  le tableau 1 avait ouverte.
%
%  ET LE MUSEE S'APPELLE « अजायबघर », LA MAISON DES MERVEILLES. Le
%  mot est persan — عجائب‌خانه — et il nomme le batiment par ce que le
%  visiteur y eprouve, non par ce qu'il contient. C'est la meme
%  maniere que le रोशनदान du tableau 1, qui nomme le vasistas par ce
%  qu'il fait ; sauf qu'ici ce n'est pas l'office qui donne le nom,
%  c'est l'effet.
%
%  LA POMPE A VAPEUR (67) EST « भाप के दमकल », ET LE MOT DIT L'ANNEE.
%  दमकल est le nom ordinaire du camion de pompiers en 2026 ; il est
%  fait de दम, le souffle ou la vapeur, et de कल, la machine — la
%  machine a souffle. Il a ete forge quand la pompe a incendie
%  fonctionnait a la vapeur, exactement comme sur cette planche de
%  1926, et il est reste sur l'engin qui a remplace celle-la. Le
%  tableau 6 avait montre le vocabulaire de la lumiere se datant tout
%  seul par ses emprunts ; celui-ci montre un mot qui se date par sa
%  MECANIQUE, et qui a survecu a l'objet qu'il decrivait.
%
%  Enfin l'administration de cette ville parle persan et ses murs
%  parlent la vieille langue. कचहरी le tribunal (22), दफ्तर le bureau,
%  चुंगीघर la douane (10), कोतवाल le commissaire sont venus avec le
%  Sultanat et les Moghols ; छावनी la caserne (16), चौक la place (96),
%  पुल le pont, बाजार, मुहल्ला etaient la avant. ON PEUT LIRE
%  L'HISTOIRE ADMINISTRATIVE D'UNE VILLE DANS LA PROVENANCE DE SES
%  BATIMENTS.
% ===================================================================

%%K t11-apar-1 apar
\VUsaut{2.0mm}
\VUcentre{13.2pt}{90}{तिसरका शृंखला}
\VUsaut{3.0mm}
\VUfilet{16mm}
\VUsaut{5.6mm}
\VUtitre{15.0pt}{60}{नकसा नं. ११}
\VUsaut{1.4mm}
\VUpk{11.4pt}{40}{सहर आ ओकर इमारत। --- आग।}
\VUsaut{2.2mm}

%%K t11-tit-1 sub
\VUpk{10.2pt}{40}{उपनगर।}
\VUsaut{3.0mm}

%%K t11-c1-01-1 p
१. --- हम एगो बड़का सहर में रहेनी जवन महासागर से सौ किलोमीटर दूर बा,
आ तबो ऊ एगो ऊँच दर्जा के \VUgras{बंदरगाह} \textsuperscript{(1)} बा। जवन
नदी ओकरा किनारे बहेले ऊ चार सौ मीटर चाकर बा आ बहुते सुंदर मोड़ बनावेले।

%%K t11-c1-02-1 p
२. --- सहर के जवन हिस्सा \VUgras{घाट} \textsuperscript{(2)} पर बा ऊ पूरा
जहाजी आ कारखाना वाला लागेला, आ एगो \VUgras{उपनगर} \textsuperscript{(3)}
बनावेला जवना में खाली मल्लाह आ मजूर रहेला। ई लोग \VUgras{कारखाना}
\textsuperscript{(4)} में आ \VUgras{गैस के कारखाना} \textsuperscript{(5)} में
काम करेला, जेकर ऊँच \VUgras{चिमनी} \textsuperscript{(6)} आ \VUgras{गैस के
टंकी} बहुते दूर से लउकेला।

%%K t11-c1-03-1 p
३. --- मध्यकाल में सहर किलाबंद रहे, आ ओकर परकोटा के कुछ अवसेस अबहीं ले
मिलेला, खास तौर पर ऊ \VUgras{फाटक} \textsuperscript{(8)}, जवन पुरनका
\VUgras{किलाबंद गढ़} \textsuperscript{(7)} के ह आ जेकरा दुनो ओर दू गो
\VUgras{बुर्ज} \textsuperscript{(9)} बा जवना के अबहीं \VUgras{जेलखाना} के तौर पर इस्तेमाल
कइल जाला।

%%K t11-c1-04-1 p
४. --- एह पूरा \VUgras{मुहल्ला} में, जवन बहुते पुरान लागेला, खाली एगो
इमारत कुछ नया बा : ऊ ह \VUgras{चुंगीघर} \textsuperscript{(10)}। ऊ घाट आ ओह
\VUgras{गली} \textsuperscript{(11)} के बीच बा जवन पुरनका \VUgras{नगरपालिका
भवन} \textsuperscript{(12)} तक ले जाले। ई आखिरी इमारत ओह बड़हन
\VUgras{घंटाघर} \textsuperscript{(13)} से आसानी से पहिचानल जाला जवन पुरनका
सहर पर छाइल बा।

%%K t11-c1-05-1 p
५. --- सड़क के दोसरा ओर \VUgras{ठाढ़} बा एगो इमारत जवन चुंगीघर के
नमूना पर बनल बा : ऊ ह \VUgras{सट्टा-बाजार} \textsuperscript{(14)}। ओकर एगो
अगवारी एगो छोट चउक की ओर देखत बा जहाँ \VUgras{कमिस्नरी}
\textsuperscript{(15)} बा। सचहूँ, हमार सहर राजधानी ना होखला के बादो जिला
के जरूरी मुख्यालय बा।

%%K t11-tit-2 sub
\VUsaut{5.0mm}
\VUpk{10.2pt}{40}{सहर के सबसे ऊँच हिस्सा।}
\VUsaut{3.4mm}

%%K t11-c1-06-1 p
६. --- सेना के टुकड़ी, जवन ठीक-ठाक बड़ बा, बहुते सेहतमंद बड़हन
\VUgras{छावनी} \textsuperscript{(16)} में रहेला ; ऊ \VUgras{नगर के
बगइचा} \textsuperscript{(17)} के जाली तक फइलल बा। जवन \VUgras{चौड़ा सड़क}
\textsuperscript{(19)} एह बगइचा तक ले जाले ऊ \VUgras{बचत बैंक}
\textsuperscript{(18)} के पाछे से गुजरेले आ \VUgras{बड़का गिरजाघर}
\textsuperscript{(20)} के चउक पर खतम होले।

%%K t11-c1-07-1 p
७. --- ई सब इमारत एगो छोट पहाड़ी पर सीढ़ीदार फइलल बा, जहाँ हमनी के
बड़हन \VUgras{माध्यमिक इसकूल} \textsuperscript{(21)} भी बा।

%%K t11-c1-07-2 p
अगर थोड़ा ढालू रस्ता से उतरल जाव, जवन बड़का सड़क से मिलेला, त
\VUgras{कचहरी} \textsuperscript{(22)} लगे पहुँचल जाला, जेकरा सामने एगो
सुहावन \VUgras{सरबजनिक बगइचा} \textsuperscript{(26)} फइलल बा। ई
\VUgras{बगइचा-चउक} बहुते बड़ नइखे, बाकिर ई सहर के बीच में हरियरी आ
ठंढक के मरूद्यान बना देला। एहीसे इहाँ सहर के लोग के बहुते आवाजाही
रहेला, जे \VUgras{कॉफीघर} \textsuperscript{(25)} के छाजन के नीचे बइठ के
ओह फौजी बजनियाँ लोग के सुनल पसंद करेला जे बियफे आ इतवार के ओह
\VUgras{मंडप} \textsuperscript{(27)} में बजावेला जवन घास के मैदान आ
फुलवारी के बीच बनल बा।

%%K t11-c1-08-1 p
८. --- ओहमें से एगो घास के मैदान पर काँसा के \VUgras{मूरत}
\textsuperscript{(28)} ठाढ़ बा, जवन हमनी के एगो सहरी के इयाद में बनल बा,
जे आपन जनम के सहर के बड़का सेवा कइले रहलें।

%%K t11-tit-3 sub
\VUsaut{4.0mm}
\VUpk{10.2pt}{40}{सवारी।}
\VUsaut{3.4mm}

%%K t11-c1-09-1 p
९. --- अबहीं ले हमनी के सहर में बिजली के अँजोर नइखे, ओहिजा खाली
\VUgras{गैस के बत्ती} \textsuperscript{(29)} बा। बाकिर \VUgras{इहाँ के
अखबार} एह बात खातिर मुहिम चलावत बा कि अँजोर के ई तरीका सुधारल जाव,
\VUgras{जइसे} \VUgras{ट्राम के खींचे के तरीका} पहिलही से सुधारल जा
चुकल बा। \VUgras{घोड़ा} से \VUgras{खींचल} \VUgras{ऊ} पुरनका गाड़ी लगभग
पूरा तरह \VUgras{बिजली के ट्राम} \textsuperscript{(31)} से बदल गइल बा,
जवन जातरी लोग के \VUgras{सहर के एक छोर से दोसरा तक} बहुते \VUgras{सस्ता}
में जल्दी पहुँचा देला।

%%K t11-c1-10-1 p
१०. --- कचहरी लगे \VUgras{बैंक} \textsuperscript{(32)} बा जहाँ बेपार के
कागज भुनावल जाला। \VUgras{बैंक के दलाल} \textsuperscript{(33)} लोग घरे-घरे
जाके बकाया वसूलेला।

%%K t11-tit-4 sub
\VUsaut{4.0mm}
\VUpk{10.2pt}{40}{कला आ पढ़ाई।}
\VUsaut{3.4mm}

%%K t11-c1-11-1 p
११. --- जवन सुंदर इमारत लगहीं ठाढ़ बा ऊ \VUgras{अजायबघर} ह। ओहमें एके
साथे सहर के \VUgras{पुस्तकालय} \textsuperscript{(34)}, \VUgras{प्रकृति-बिज्ञान
के सुंदर संग्रह} \textsuperscript{(35)} आ एगो सुघर \VUgras{चित्रसाला}
\textsuperscript{(36)} बा, जवन एगो शानदार पक्का \VUgras{प्रदर्सनी}
बनावेला।

%%K t11-c1-11-2 p
ऊ हफ्ता में दू बेर जनता खातिर खुलेला, बाकिर परदेसी लोग कवनो भी दिन
\VUgras{रखवार} \textsuperscript{(37)} के थोड़ा बकसीस देके ओकरा के देख सकेला।
\VUgras{चित्रकार} \textsuperscript{(38)} लोग ओहिजा काम करे आवेला। ओकनी के
आपन तिपाई के \VUgras{सामने}, हाथ में \VUgras{रंग-पटरी} लेके, नामी चित्र
के नकल करत देखल जाला।

%%K t11-c1-12-1 p
१२. --- ऊँचाई पर, अजायबघर के पाछे, ऊ \VUgras{गुंबद}
\textsuperscript{(40)} लउके लागेला जवन \VUgras{अस्पताल}
\textsuperscript{(39)} के ह, आ ओह
\VUgras{मास्टर के इसकूल} \textsuperscript{(41)} के इमारत भी, जहाँ हमनी
के नगर के इसकूलन के मास्टर लोग के तालीम मिलत रहे।

नाटकघर के चउक पर एगो \VUgras{निजी घर} \textsuperscript{(43)} में
\VUgras{डाकघर} \textsuperscript{(42)} बा।

%%K t11-tit-5 sub
\VUsaut{5.0mm}
\VUpk{11.4pt}{40}{आग।}
\VUsaut{3.0mm}

%%K t11-tit-6 sub
\VUpk{10.2pt}{40}{आग के हल्ला।}
\VUsaut{3.4mm}

%%K t11-c2-01-1 p
१. --- सहर में एगो डरावन खबर फइल गइल : नाटकघर में अभिए आग लाग गइल !
जवन बेरा हम \VUgras{चउक} \textsuperscript{(96)} पर पहुँचनी, भीड़
पहिलही से \VUgras{पटरी} \textsuperscript{(46)} पर आ \VUgras{सड़क}
\textsuperscript{(47)} पर जुट चुकल रहे। \VUgras{डाकिया}
\textsuperscript{(44)} आ \VUgras{चिट्ठी बाँटे वाला} \textsuperscript{(45)}
डाकघर से बाहर निकलत बाड़ें। एगो \VUgras{ट्राम चलावे वाला}
\textsuperscript{(48)} के आपन गाड़ी रोके के पड़ल ताकि सहर के
\VUgras{एंबुलेंस} \textsuperscript{(50)} गुजर सके, जवन एगो \VUgras{जर्राह}
\textsuperscript{(52)} आ एगो \VUgras{नर्स} \textsuperscript{(51)} के लेके
आवत बा एगो \VUgras{घायल} \textsuperscript{(53)} के देखे खातिर, जेकरा के
\VUgras{नक्कालात} \textsuperscript{(54)} पर \VUgras{अखबार के गुमटी}
\textsuperscript{(55)} लगे ले आवल गइल बा।

%%K t11-c2-02-1 p
२. --- पैदल सेना के एगो टुकड़ी, जे अभ्यास से लवटत रहे, रुक गइल ताकि
\VUgras{सहर के घोड़सवार पुलिस} \textsuperscript{(61)} के साथे मिल के
बेवस्था सँभाल सके। \VUgras{घोड़सवार लोग}, जेकर घोड़ा अगुड़ी उठावत बा,
\VUgras{फौजी सिपाही} \textsuperscript{(57)} भा \VUgras{जेंडार्म}
\textsuperscript{(63)} से जादे नीक से \VUgras{तमासबीन} \textsuperscript{(56)}
के पाछे हटावत बाड़ें। ऊपर से जेंडार्म लोग बहुते बेस्त बा। ओहमें से एगो
\VUgras{पुलिस वाला} \textsuperscript{(64)} के बतावत बा कि दोसरका
\VUgras{दुर्घटना में घायल} \textsuperscript{(65)} के कहाँ ले जाइल जाव —
एगो हिम्मती दमकल वाला जे सीढ़ी से गिर गइल बा।

%%K t11-c2-03-1 p
३. --- \VUgras{अफसर} \textsuperscript{(66)} बतावत बाड़ें कि आवत
\VUgras{भाप के दमकल} \textsuperscript{(67)} के कहाँ लगावल जाव। ओह ताकतवर
मसीन के अगोरत ऊ फुरती से एगो \VUgras{हाथ के पंप} \textsuperscript{(68)}
लगववले बाड़ें, जेकर लमहर \VUgras{नली} \textsuperscript{(69)} आग पर धार
नियर पानी फेंकत बा।

छव गो दमकल वाला ओकरा के चलावत बा, जबकि ओह लोग के संगी, जे
\VUgras{नोजल} \textsuperscript{(70)} थामले बा, \VUgras{धार}
\textsuperscript{(71)} के आग के बीचोबीच ले जात बा।

%%K t11-c2-04-1 p
४. --- नाटकघर के दोसरा ओर बड़का \VUgras{सीढ़ी} \textsuperscript{(72)}
टेकावल गइल बा। ओहसे दमकल वाला लोग ओह लपट लगे पहुँच पावत बा जवन करिया
\VUgras{धुँआ} \textsuperscript{(75)} के बवंडर के बीच से निकलत बा।

%%K t11-tit-7 sub
\VUsaut{5.0mm}
\VUpk{10.2pt}{40}{बहादुरी के काम।}
\VUsaut{3.4mm}

%%K t11-c2-05-1 p
५. --- \VUgras{आग} \textsuperscript{(74)} \VUgras{नाटकघर}
\textsuperscript{(73)} के बरामदा में लागल, जब ऊ \VUgras{ओपेरा} के रिहर्सल
होत रहे जेकर पहिलका \VUgras{खेल} काल्ह होखे वाला रहे। गनीमत बा कि
दरसक-कमरा में थोड़े \VUgras{दरसक} \textsuperscript{(76)} रहलें। बेचारा लोग
\VUgras{बालकनी} \textsuperscript{(77)} पर सरन लेलस, जहाँ हिम्मती
\VUgras{दमकल वाला} \textsuperscript{(86)} लोग सीढ़ी आ रस्सा से ओकनी के
बचावे आवत बाड़ें।

%%K t11-c2-06-1 p
६. --- \VUgras{खंभा वाला बरामदा} \textsuperscript{(78)} के नीचे, जहाँ
\VUgras{इस्तिहार} \textsuperscript{(79)} खेल के खबर देत बा, ओहिजा
\VUgras{बजनियाँ} \textsuperscript{(81)} लोग के \VUgras{आरकेस्ट्रा}
\textsuperscript{(80)} से भागत देखल जात बा, आपन बाजा लेके :
\VUgras{बेला} \textsuperscript{(81)}, \VUgras{बाँसुरी} \textsuperscript{(82)},
\VUgras{चेलो} \textsuperscript{(83)}, \VUgras{ट्रंबोन} \textsuperscript{(84)},
\VUgras{ढोल} \textsuperscript{(85)}, वगैरह। लोग \VUgras{सामान}
\textsuperscript{(87)} के कुछ हिस्सा बचावे के कोसिस करत बा।

%%K t11-c2-07-1 p
७. --- \VUgras{नट} \textsuperscript{(89)} आ \VUgras{नटी}
\textsuperscript{(88)}, \VUgras{गवइया} आ \VUgras{गवइनी} के आपन नाटक के
\VUgras{पोसाक} \textsuperscript{(90)} में भागे के पड़ल। तेनोर एगो
\VUgras{पत्रकार} \textsuperscript{(92)} के बतावत बाड़ें कि ई सब कइसे भइल।
\VUgras{रिपोर्टर} पहिलही पल से अफसर लोग के साथे दउड़ के आ गइल रहलें :
\VUgras{कमिस्नर} \textsuperscript{(91)}, \VUgras{नगर-मुखिया}
\textsuperscript{(93)}, \VUgras{पुलिस कप्तान} \textsuperscript{(94)} आ
\VUgras{अदालत के अफसर} \textsuperscript{(95)}, जे जाँच करिहें ताकि तय
होखे कि केकर जिम्मेवारी बा।
